\section{Background and Method}
\label{p11:sec:background}

\noindent
This section surveys the frameworks the paper draws upon, states their relation to one another, and declares the methodological concession under which the paper's assessments are made. The frameworks are surveyed not as competitors among which the paper selects, but as the resources the cycle of a shared life draws upon at the moments suited to each. Five are distinguished, by the epistemology of flourishing each carries.

\subsection{Five frameworks for the cognition and practice of flourishing}
\label{p11:sec:bg-five}

\emph{Positivism}, in the sense the paper employs, is the position that flourishing is an observable and measurable matter of fact, accessible through its indicators and tractable by quantification. Its developed form in the study of well-being is the tradition of subjective well-being, in which well-being is operationalised as life satisfaction together with the balance of positive over negative affect, and measured by validated instruments \parencite{diener1985}. The hedonic tradition that underwrites these measures is gathered in \textcite{kahneman1999}, and its instruments include scales of positive and negative affect \parencite{watson1988} and of quality of life \parencite{whoqol1998}. Alongside the subjective measures stand the objective indicators of quality of life, the measurable conditions of a life, and the apparatus of psychometrics by which any such measure is assessed for reliability and validity. The positivist commitment is that what cannot be brought, however indirectly, to measurement cannot be a matter of knowledge in the strict sense.

\emph{Evidentialism}, as a practical epistemology, is the position that a belief, and the practice founded on it, ought to be proportioned to the evidence, and revised as the evidence revises. In the study of flourishing its developed form is evidence-based practice, in which an intervention is warranted by the weight of controlled evidence that it produces the effect claimed for it. Evidentialism is not identical to positivism, though the two are allied: positivism concerns what flourishing is, an observable fact, while evidentialism concerns how belief about flourishing ought to be formed and updated, by proportioning to evidence. The distinction matters to the paper, because evidentialism, in its deeper form, is the epistemology of belief revision as such, and so reaches below the explicit weighing of studies to the continuous updating of belief that the paper locates at the basis of all appraisal (\xref{p11:sec:appraisal}).

\emph{Normative theory} is the position that flourishing is not exhausted by what is the case but involves what ought to be valued: that a good life is constituted by ends there is reason to hold worthy, and not merely by the satisfaction reported or the conditions observed. Its most developed contemporary form, and the one most consonant with this series, is the capabilities approach, in which the measure of a life is the substantive freedom of a person to be and to do what that person has reason to value, the capabilities rather than the achieved functionings alone \parencite{sen1999,nussbaum2011}. The approach has a precedent in the recovery of an Aristotelian conception of flourishing as vulnerable to circumstance \parencite{nussbaum2001}. The capabilities approach is consonant with the series because its central notion, the substantive freedom to unfold what one has reason to value, is closely allied to the conception of flourishing as the unfolding of each party's own dynamics that the prior papers developed.

\emph{Virtue ethics} is the position that flourishing, eudaimonia in its classical sense, is the activity of a life in accordance with excellence of character, and that such excellence is acquired through habituation. In Aristotle the virtues are states of character, formed by the repeated performance of the corresponding actions until they become a settled disposition, a \emph{hexis}; eudaimonia is then the activity of the soul in accordance with these acquired excellences over a complete life \parencite{aristotle2009}. The contemporary recovery of this conception, against the modern reduction of ethics to rules, is owed in large part to \textcite{macintyre1981}, and its extension to the virtues of acknowledged dependence to \textcite{macintyre1999}. The Confucian tradition of self-cultivation supplies a parallel account, in which the person is formed through sustained cultivation, ritual practice, and the ordering of the household and its manners, the good life being an attained and maintained cultivation rather than a state secured once \parencite{tu1985}. The classical sources of this tradition include the \emph{Analects} \parencite{confucius_analects} and the doctrine of the mean \parencite{zhongyong}. Both traditions locate flourishing in the formation of disposition through habit, and both thereby connect the cognition of flourishing to its practice in a way the other frameworks do not, which is why virtue ethics is the framework proper to the practical moment of the cycle (\xref{p11:sec:practice}).

\emph{Phenomenology and hermeneutics} together name the position that flourishing is lived and interpreted before it is measured: that it is given, in the first instance, as the felt quality and the understood meaning of a life, and that this givenness is not a defective approximation to a measurable fact but the matter itself, to which any measure is answerable. The phenomenological tradition supplies the account of the lived and intersubjective texture of experience on which the prior papers drew; the account of perception and the lived body is owed to \textcite{merleau1945}, and the diagnosis of the reduction of the lived world to its measurable indicators to \textcite{husserl1936}; the hermeneutic tradition supplies the account of understanding as the interpretation of meaning within a circle that cannot be exited for a view from nowhere. Against the positivist commitment, this framework holds that the reduction of flourishing to its measurable indicators discards precisely what flourishing is, and that the understanding of a shared life is closer to the interpretation of a text than to the reading of an instrument. The structure of such understanding, as a circle that cannot be exited for a view from nowhere, is developed by \textcite{gadamer1960}.

\subsection{The frameworks as symmetry-broken phases}
\label{p11:sec:bg-phases}

These five frameworks are not five candidate theories of one measurable object, among which evidence will eventually decide. They are five phases, in the sense the prior paper gave the term, of the cognition and practice of flourishing: each is the determinate form the understanding of flourishing assumes under the conditions of a particular discipline and a particular task, and none is the undivided structure of which the others are partial views. The positivist measures what can be measured and is silent on what cannot; the phenomenologist describes what is lived and resists its reduction to measure; the normative theorist specifies what ought to be valued and cannot guarantee its realisation; the virtue theorist describes the formation of disposition and cannot legislate character; the evidentialist proportions belief to evidence and is mute where evidence cannot reach. Each says, in its own broken phase, what no other can say, and each is silent where another speaks.

The relation among them is therefore not one to be resolved by the victory of one framework, nor by their eclectic combination into a single super-framework, which would be the metalanguage the series has argued does not exist. It is the relation of a polyphony to a truth that no single voice possesses, and the burden of the cycle is to draw upon each at the moment to which it is suited, without subordinating the others to it. This is the methodological inheritance of the prior paper, transposed from the disciplines of the theoretical papers to the frameworks of the practical one, and it is a dialectical-materialist rather than a relativist commitment: the frameworks are partial determinations, each from its position, of a flourishing that is real and in production, and not so many incommensurable language games among which there is no fact.

\subsection{The operational concession}
\label{p11:sec:bg-concession}

A practice requires what is operable and what is assessable. The cultivation of a shared life proceeds by actions that can be taken and revised, and its revision requires some assessment of whether the cultivation tends toward the good. But the prior papers established that the value at issue in a relation is transversal across the three registers, emergent on their coupled dynamics, and not exhausted by any symbolisation; and they established that to fix such a value in a determinate, possessable structure, to cash it, is the operation of settling, and the occupation of the place of the barred Other. There is, accordingly, a tension at the centre of any practical and assessing treatment of flourishing: assessment requires that value be structured into an assessable form, and the value that is real resists exactly that structuring.

The paper meets this tension by a concession it states openly and bears throughout. Where it must assess, it restricts itself to value as structured into an assessable proxy, and it holds that this proxy is not the value the prior papers established to be irreducible. The structured proxy is an instrument of practice, warranted by its operational use, and it is not a representation of the real value, which it cannot be. This is the first occasion in the series on which an incomplete value is adopted deliberately rather than diagnosed as a failure, and the adoption is the honest discharge of the prior result rather than its reversal. Since the real value cannot be cashed without occupying the place of the Other, the honest course, where a practice must assess, is not to claim that the assessment reaches the real value, which would be that very occupation, but to concede that what is assessed is a proxy, to use it for what it is operationally worth, and to bear responsibility for its incompleteness. The concession governs every assessment the paper makes, and most consequentially the quantitative assessment of justice in \xref{p11:sec:evaluation}, whose results are to be read as operationally valid and not as the measurement of a relation's real worth.

The concession also fixes the standing of the explicit, instrumented evaluation that the positivist framework supplies. That evaluation is not thereby disparaged. It is, on the contrary, indispensable to a practice, since a practice without assessment cannot be revised; and the structured proxy it yields is, within its operational limits, a real gain, public, comparable, and accumulable in a way that the unstructured appraisal is not. The concession is not a verdict against measurement but a specification of its standing: the measure is an operationally warranted proxy for a value it does not reach, to be used as such and not mistaken for the value itself. The paper now turns to the vocabulary in which the cycle and its regimes will be described.
